\section{Positive matrices and attained lower energies}

All matrices in this section act on a finite-dimensional complex Euclidean
space. The order is the positive-semidefinite order.

\begin{theorem}[Energies of normalized eigenvectors]
    \label{thm:normalized_eigenvector_energies}
    \lean{Matrix.PosSemidef.nonneg_of_normalized_eigenvector,
        Matrix.groundEnergy_eq_of_normalized_eigenvectors}
    \leanok
    Let $Ax=ex$, where $e\in\R$ and $\|x\|=1$. If $A\ge0$, then $e\ge0$.
    If also $Ay=e'y$ with $e'\in\R$, $\|y\|=1$, and $A\ge e\Id$, $A\ge e'\Id$, then $e=e'$.
    The latter assertion does not require $A\ge0$.
\end{theorem}
\begin{proof}
    \leanok
    Positivity gives $e=\langle x,Ax\rangle\ge0$. In the second assertion,
    apply this observation to $(A-e\Id)y=(e'-e)y$ and
    $(A-e'\Id)x=(e-e')x$ to obtain both inequalities.
\end{proof}

\begin{theorem}[A common extremal vector in an order sandwich]
    \label{thm:order_sandwich_eigenvector}
    \lean{Matrix.mulVec_eq_smul_of_order_sandwich}
    \leanok
    Suppose $e\Id\le F\le A$ and $Ax=ex$, where $e\in\R$.
    Then $Fx=ex$; the vector need not be normalized or nonzero.
\end{theorem}
\begin{proof}
    \leanok
    The positive matrices $F-e\Id$ and $A-F$ have a sum annihilating $x$.
    The sum of their nonnegative quadratic forms at $x$ is zero, so each
    annihilates $x$.
\end{proof}

\begin{theorem}[Ground energy above a positive connecting interaction]
    \label{thm:positive_boundary_ground_energy}
    \lean{Matrix.groundEnergy_bounds_of_positive_boundary_and_zero_test}
    \leanok
    Let $X,Y\ge0$, $Y\le C\Id$, and let $x,z$ be unit vectors. Suppose
    $Xx=0$, $(X+Y)z=ez$, and $X+Y\ge e\Id$, where $C,e\in\R$.
    Then $0\le e\le C$. These are the energy bounds used
    in~\cite[lines 873--886 and 1010--1024]{Arad2016RigorousRG}.
\end{theorem}
\begin{proof}
    \leanok
    \uses{thm:normalized_eigenvector_energies}
    Positivity at $z$ gives $e\ge0$. The lower bound at $x$ gives
    $e\le\langle x,(X+Y)x\rangle=\langle x,Yx\rangle\le C$.
\end{proof}

\begin{theorem}[A finite sum of positive contractions]
    \label{thm:positive_unit_sum_budget}
    \lean{Matrix.sum_posSemidef_and_le_card_of_positive_unit_terms}
    \leanok
    For a finite set $S$ and matrices satisfying $0\le h_i\le\Id$ for
    $i\in S$, the matrix $\sum_{i\in S}h_i$ is positive and at most $|S|\Id$.
    The matrix index set and $S$ may be empty.
\end{theorem}
\begin{proof}
    \leanok
    Sum the lower and upper operator inequalities over $S$.
\end{proof}

\begin{theorem}[Centering the original spectral basis]
    \label{thm:original_basis_centered_gap}
    \lean{Matrix.IsHermitian.original_basis_centered_gap_data}
    Let $H$ be Hermitian on a nonzero finite-dimensional complex coordinate
    space. Suppose $\mu\in\operatorname{spec}(H)$, $g>0$, and every spectral
    value is either $\mu$ or at least $\mu+g$. Let $s,E_0\in\R$, and suppose
    a unit vector $z$ satisfies
    $(H-s\Id)z=E_0z$ and $H-s\Id\ge E_0\Id$. Then $\mu=E_0+s$.
    Every member $v_j$ of the original orthonormal eigenbasis of $H$ satisfies
    \begin{align}
        (H-s\Id)v_j & =(\lambda_j-s)v_j,\notag                     \\
        \lambda_j-s & =E_0\quad\text{or}\quad E_0+g\le\lambda_j-s.
    \end{align}
\end{theorem}
\begin{proof}
    \uses{thm:normalized_eigenvector_energies}
    The supplied unit vector attains the lower bound of $H-s\Id$.
    A unit $\mu$-eigenvector attains the lower bound $\mu$ of $H$.
    Comparing their Rayleigh quotients in both directions gives
    $\mu=E_0+s$. Subtract $s$ in each original eigenvector equation and in
    its spectral alternative. This retains the original basis, as needed
    for the comparison in~\cite[lines 909--960]{Arad2016RigorousRG}.
\end{proof}
